\section{Requirements not yet discharged}
This part records unverified source requirements. None of its unavailable
original-source claims is used as a proved premise in Part I. The report
remains marked \texttt{INCOMPLETE} while these requirements are outstanding.

\subsection{B.3, subject 21: the two prescribed original articles}
The affine henselian rigidity statement and the completion-complement
comparison have complete checked deductions in
Theorems~\ref{thm:B3-21} and~\ref{thm:B3-fujiwara}, from the opened
Stacks, Haine--Holzschuh--Wolf and Bhatt--Mathew sources. In particular,
the latter original paper's noetherian hypothesis is distinguished
from the subsequent nonnoetherian generalization.

The separately requested direct original readings remain unavailable:
Gabber, \emph{Affine analog of the proper base change theorem},
Israel Journal of Mathematics 87 (1994), pp.~325--335,
\S5, Corollary 1; and Fujiwara, \emph{Theory of tubular neighborhood in
\'etale topology}, Duke Mathematical Journal 80 (1995), pp.~15--57,
Corollary 6.6.4. These theorem locators have been checked in the opened
later sources, not in the originals themselves
\cite{BGabber1994,BFujiwara1995}.

The retrieval record includes publisher and historical publisher
endpoints, author and institutional searches, DOI and repository
metadata, public web-archive listings and a public archived landing
page, and whole-journal catalogue records. The exact archived Duke
issue marks its PDF and OCR private; no private file was requested.
The archived article landing page denies full-text access, whereas
the current non-PDF response is an anti-bot page. We do not conflate
these distinct failures. No challenge or subscription control was
bypassed. The smallest remaining gate is direct access to the two
specified original passages, followed by a hypothesis-by-hypothesis
comparison with the formulations already proved here.

\subsection{Scope boundaries which are not claimed as solved problems}
The raw inventory of the twenty-eight thesis subjects and nine
\emph{Pour m\'emoire} headings records the strongest specified proved
formulation used here. Broad historical headings are not silently
converted into blanket theorems. In particular, we do not assert an
unrestricted abelian category of mixed motives, resolution in arbitrary
positive characteristic, every question in SGA 2, or a formal meaning
of the undefined category of ``sch\'emas morcel\'es''. The corresponding
positive results, counterexamples and exact scope are stated in Part B.
This is the explicit ``strongest proved statement'' alternative in
task B.3, rather than a claim that the whole historical programme has
been affirmatively settled.

The positive-characteristic all-prime abelian-fourfold conclusion
explicitly uses Broe's 2026 preprint in
Theorem~\ref{thm:B1b-abelian}. We checked the theorem statement and the
deduction required here, and do not claim an independent verification
of the entire preprint. Ancona's published numerical positivity alone
is not represented as proving that conclusion.

The class requested in B.1(b) is the precisely stated domain of the
componentwise-definiteness argument: algebraic Lefschetz projections
and positivity on the homological primitive algebraic classes.
The criterion in Proposition~\ref{prop:B1b-criterion} proves the
deduction on that entire domain. Its hypotheses are separately
verified for the classes identified in Part B. We do not replace
positivity on that homological space by positivity merely on its
numerical quotient. A geometric classification of all isomorphism
types satisfying the criterion is not asserted or introduced as an
additional requirement of the prompt.

\subsection{The remaining named-source delivery requirement}
The prompt separately instructs us to fetch its named historical works.
This is an outstanding delivery requirement even where an independent
proof or an opened replacement proves the mathematics. In particular,
the source audit has identified outstanding original-text readings for
Lehto--Virtanen, Pfluger, Alling--Greenleaf and Lieberman. The requested Raynaud
thesis is \emph{Faisceaux amples sur les sch\'emas en groupes et les
espaces homog\`enes}, LNM119 (1970), distinct from the opened Picard
article. The Birman book's relevant mapping-class chapter remains
outside the original preview that was opened.

The originals of Nakamura (1990), Ihara--Nakamura (1997), and
Gabriel (1932) have now been obtained and the relevant mathematical
passages checked. Their exact hypotheses and their limits are recorded
in the corresponding parts and bibliography. These three source
requirements are closed.

This list is a live snapshot: the report and bibliography record which
originals have subsequently been obtained, the exact passages actually
read, and the remaining inaccessible works. Original SGA4, $4\frac12$,
5 and6 texts have now been retrieved and selected exact passages read.
Lavrentiev's original 1934 text and Bers's original 1974 announcement
have been opened; the bordered conventions are checked in
\cite[\S\S2(F),3(F),4(E)]{ABers1981}. An announcement of a theorem is
not treated as inspection of an unpublished proof. Koebe's original
1907 covering-radius passage is now read in
\cite[pp.~203--205]{AKoebe1907}. The joint Earle--Marden original
author draft dated24August2011 is now opened and its local construction
derived in Proposition~\ref{A1-Earle-Marden-original}; this closes
A.1(c)'s joint-source requirement. The separate2012 chapter remains
unopened, but the prompt did not specify that chapter as the sole
acceptable joint source. We do not call these
named originals optional or silently replace
them by later accounts. For each mathematical conclusion used in Part I,
the proof instead identifies its actual opened source or supplies the
argument directly.
